\documentclass[24pt,letterpaper]{article}
\usepackage[left=1cm,right=1cm,top=1cm,bottom=1cm]{geometry}
\usepackage{graphicx} % Required for inserting images
\usepackage{parskip}

\title{Hoja Referencial MIPS 32}
\author{Maivet Pereira}
\date{May 2026}

\begin{document}

\maketitle

\section{Aritmetica}
add(suma) \hspace{0.60cm} add \$s1,\$s2,\$s3  \hspace{0,60cm} \# \$s1 = \$s2 + \$s3

subtract(resta) \hspace{0.60cm} sub \$s1,\$s2,\$s3  \hspace{0,60cm} \# \$s1 = \$s2 - \$s3

add immediate \hspace{0.60cm} addi \$s1,\$s2,100  \hspace{0.60cm} \# \$s1 = \$s2 + 100

\section{Transferencia de Datos}
w = Word (Palabra entera = 4 bytes = El tamaño de un int normal).

h = Halfword (Media palabra = 2 bytes = El tamaño de un short).

b = Byte (Un solo byte = 1 byte = El tamaño de un char).

------------------------------------------------------------------------------------------------------------------------------------------------------------------

load word \hspace{0.60cm} lw \$s1, 4(\$s2)\$s1 = TRAE un entero desde la RAM (dirección \$s2 + 100)

store word \hspace{0.60cm} sw \$s1, 8(\$s2) \hspace{0.60cm} \# GUARDA el entero de \$s1 en la RAM (dirección \$s2 + 100)

load halfword \hspace{0.60cm} lh \$s1, 100(\$s2) \hspace{0.60cm} \# \$s1 = TRAE medio entero  desde la RAM (dirección \$s2 + 100)

store halfword \hspace{0.60cm} sh \$s1, 100(\$s2) \hspace{0.60cm} \# GUARDA solo la mitad de \$s1 en la RAM (dirección \$s2 + 100)

load byte \hspace{0.60cm} lb \$s1, 100(\$s2) \hspace{0.60cm} \# \$s1 = TRAE un solo byte desde la RAM (dirección \$s2 + 100)

store byte \hspace{0.60cm} sb \$s1, 100(\$s2) \hspace{0.60cm} \# GUARDA solo el último byte de \$s1 en la RAM (dirección \$s2 + 100)

load upper immediate \hspace{0.60cm} lui \$s1, 100 \hspace{0.60cm} \# Mete el número 100 en la parte ALTA (izquierda) de \$s1 (y llena de ceros el resto)

\section{Logica}
and \hspace{0.60cm} and \$s1, \$s2, \$s3 \hspace{0.60cm} \# \$s1 = \$s2 Y \$s3 (bit a bit, da 1 solo si ambos bits son 1)

or \hspace{0.60cm} or \$s1, \$s2, \$s3 \hspace{0.60cm} \# \$s1 = \$s2 O \$s3 (bit a bit, da 1 si al menos uno de los bits es 1)

xor \hspace{0.60cm} xor \$s1, \$s2, \$s3 \hspace{0.60cm} \# \$s1 = \$s2 o exclusivo \$s3 (da 1 si los bits son diferentes)

nor \hspace{0.60cm} nor \$s1, \$s2, \$s3 \hspace{0.60cm} \# \$s1 = NO (\$s2 O \$s3) (invierte el resultado del O, util para hacer un NOT)

and immediate \hspace{0.60cm} andi \$s1, \$s2, 100 \hspace{0.60cm} \# \$s1 = \$s2 Y el número 100 (para comparar un registro con una constante)

or immediate \hspace{0.60cm} ori \$s1, \$s2, 100 \hspace{0.60cm} \# \$s1 = \$s2 O el número 100 (se usa mucho junto a lui para cargar constantes grandes)

shift left logical \hspace{0.60cm} sll \$s1, \$s2, 2 \hspace{0.60cm} \# \$s1 = Desplaza los bits de \$s2 a la IZQUIERDA 2 veces (Equivale a multiplicar por 4)

shift right logical \hspace{0.60cm} srl \$s1, \$s2, 2 \hspace{0.60cm} \# \$s1 = desplaza los bits de \$s2 a la DERECHA 2 veces (equivale a dividir entre 4)

\section{Salto Condicional}
branch if equal \hspace{0.60cm} beq \$s1, \$s2, label \hspace{0.60cm} \# si (\$s1 == \$s2) salta a la etiqueta label

branch if not equal \hspace{0.60cm} bne \$s1, \$s2, label \hspace{0.60cm} \# si (\$s1 != \$s2) salta a la etiqueta label

\section{Salto Incondicional}
jump \hspace{0.60cm} j label \hspace{0.60cm} \# salto incondicional (va directo a la etiqueta label sin preguntar)

jump and link \hspace{0.60cm} jal funcion \hspace{0.60cm} \# llama a una funcion y guarda la direccion de retorno en \$ra

jump register \hspace{0.60cm} jr \$ra \hspace{0.60cm} \# salta a la direccion guardada en \$ra (se usa para volver de una funcion)

set on less than \hspace{0.60cm} slt \$s1, \$s2, \$s3 \hspace{0.60cm} \# si (\$s2 < \$s3) \$s1 = 1; si no \$s1 = 0 (sirve para los if de menor que)

set on less than immediate \hspace{0.60cm} slti \$s1, \$s2, 100 \hspace{0.60cm} \# si (\$s2 < 100) \$s1 = 1; si no \$s1 = 0 (compara un registro con un numero)

\section{Formato Basico de Instruccion}
\vspace{0.4cm}

\vspace{0.2cm}
\noindent \textbf{R} \\
\begin{tabular}{|c|c|c|c|c|c|}
\hline
cod oper (6 bits) & rs (5 bits) & rt (5 bits) & rd (5 bits) & desplaz (5 bits) & func (6 bits) \\ \hline
\end{tabular}

\vspace{0.3cm}
\noindent \textbf{I} \\
\begin{tabular}{|c|c|c|c|}
\hline
cod oper (6 bits) & rs (5 bits) & rt (5 bits) & inmediato (16 bits) \\ \hline
\end{tabular}

\vspace{0.3cm}
\noindent \textbf{J} \\
\begin{tabular}{|c|c|}
\hline
cod oper (6 bits) & direccion (26 bits) \\ \hline
\end{tabular}
\end{document}
